\section{Discussion}
\label{sec:discussion}

\subsection{What is fixed by the continuum theory fragment}

At the continuum level (treated here as the long-wavelength effective limit of a cubic lattice substrate), the theory specifies:
(i) the substrate dynamics via an isotropic hyperelastic action \eqref{eq:action},
(ii) the existence of an ordered narrowband sector that supports a controlled multiscale envelope description,
and (iii) the geometric connection structure (Berry/Wilczek--Zee) induced by polarization transport.
Given these commitments, gauge degrees of freedom are not introduced as independent fundamental fields; they arise as
connection data of the ordered carrier sector.

\subsection{Minimal-ingredient reconstruction and the role of effective theories}

The guiding stance is that established field theories may be \emph{effective} descriptions of stable, repeatable
excitation sectors of a deeper substrate.
In that reading, the Standard Model is not ``wrong''; it is a remarkably accurate macroscopic language for the
degrees of freedom that remain after coarse-graining.
The present framework aims to reduce the number of primitive ingredients by shifting explanatory weight from
``many fundamental objects'' to ``one object with richer nonlinear and topological behavior''.
Concretely, this means:
relativistic kinematics is tied to the dominant wave cone of a linear sector,
gauge structure is tied to holonomy of an ordered sector,
and particle-like behavior is tied to localized nonlinear modes.

\subsection{Lattice substrate: ontology, anisotropy, and continuum recovery}
\label{subsec:discussion-lattice}

The cubic lattice is treated not merely as a numerical convenience but as a candidate \emph{microphysical}
substrate.
In this reading, the isotropic continuum action serves as a compact \emph{effective} description valid at wavelengths
$\lambda\gg a$ (with $a$ the lattice spacing), while the lattice determines the high-$k$ dispersion and the available
polarization/mixing structure.

Because cubic symmetry is weaker than full rotational symmetry, the lattice generically produces
direction-dependent wave speeds, dispersion, and branch splitting (``birefringence'') at momenta approaching the
Brillouin-zone boundary.
Two complementary requirements follow:
(i) \emph{observational viability} demands that anisotropies at accessible scales be suppressed, and
(ii) the theory should explain \emph{why} effective isotropy emerges robustly for composite measuring devices that are
themselves excitations of the substrate.
Operationally, the paper commits to the controllable regime where the continuum limit is well-defined and
anisotropies are parametrically small; the lattice picture then becomes a concrete hypothesis about the UV completion.

\subsection{Axis-aligned triplet phases and the coupling-driven transition to non-Abelian holonomy}
\label{subsec:discussion-axis-triplet}

The ``axis-aligned Berry phase'' idea is best understood as a \emph{basis-dependent} description of a transported
near-degenerate carrier sector.
On a cubic lattice, it is natural to organize a narrowband carrier triplet
$\Psi=(\psi_x,\psi_y,\psi_z)^\top$ aligned with the lattice axes.
If inter-axis mixing is negligible ($g_{\mathrm{mix}}\to 0$), each component admits an independent rank-1 Berry connection and one
may speak meaningfully of three separate $U(1)$ geometric phases.

For any nonzero mixing ($g_{\mathrm{mix}}>0$), however, the physical eigenmodes are mixtures of axis labels.
In that regime, ``the Berry phase of the $x$ axis'' is not gauge-invariant; the invariant content is the
Wilczek--Zee holonomy of the transported subspace and its curvature.
This provides a sharp criterion for when axis-wise phase bookkeeping is legitimate (effectively decoupled channels)
and when the description must be upgraded to a non-Abelian $U(3)$ connection with an $SU(3)$ traceless component.

\subsection{Contact with lattice gauge theory and lattice-QCD Berry curvature diagnostics}
\label{subsec:discussion-lqcd}

Standard lattice QCD uses a lattice primarily as a \emph{regulator} for a continuum gauge theory; in contrast, this
framework treats the lattice as \emph{ontic} and interprets gauge structure as emergent connection data of an
ordered carrier sector.
Despite this difference in interpretation, the \emph{diagnostic machinery} transfers almost verbatim:
Berry connections and curvatures can be computed on a discretized Brillouin zone using link variables and plaquettes,
in close analogy to lattice-QCD Berry-curvature constructions \citep{PuYamamoto2018BerryLatticeQCD,Yamamoto2016BerryPhaseLatticeQCD,Fukui2005ChernDiscretizedBZ}.
A concrete gauge-invariant extraction protocol and artifact controls are summarized in \Cref{subsec:diagnostic-protocol}.

This analogy clarifies a central structural point about the ``color'' direction.
Non-Abelian Berry curvature requires (approximate) degeneracy of the transported subspace; interactions or mixing that
lift degeneracy can reduce the effective geometry to an Abelian sector.
In this model, the mixing factor $g_{\mathrm{mix}}$ plays an analogous role: it controls both the mixing of the triplet and
the gap structure that determines whether a non-Abelian adiabatic transport picture is justified.
A practical program suggested by this link is:
(i) compute $U(1)$ and $U(3)$ Berry curvatures on the momentum lattice for the relevant branches,
(ii) compare holonomies under loops in $\mathbf{k}$-space to ``charge'' labels of long-range interactions, and
(iii) test whether the traceless part exhibits genuinely non-Abelian features (e.g.\ noncommuting holonomies) in the
parameter regime where localized modes and an ordered narrowband sector coexist.

\subsection{Gravity as contraction from transverse excitation}

The extra embedding dimension is not introduced only to host additional polarizations; it also enforces a specific
geometric backreaction: transverse variation changes intrinsic distances (\Cref{subsec:gravity-channel}).
In the discrete picture this is the Pythagorean spring-length effect, and in the continuum it is already visible in
$g_{ij}=h_\star^2\,\delta_{ij}+\partial_i u\,\partial_j u$.
Under homogeneous pre-tension, this backreaction produces an inward lateral pull (contraction) around localized
transverse excitation.
In the intended weak-field regime, the ``gravitational field'' is therefore identified with a slowly varying,
coarse-grained strain/dilation pattern sourced by excitation energy (especially transverse gradient energy), and
``gravitational mass'' is tied to that stored elastic energy.

What remains open is the quantitative reduction: one would like to derive a Newtonian limit in which an effective
scalar potential built from coarse-grained strain or energy density obeys a Poisson-type equation with a universal
coupling constant, and then relate that constant to $G$.
On the simulation side, the central check is conceptually simple: create a localized transverse excitation, measure
the induced steady contraction field in the surrounding brane, and quantify how it focuses or accelerates long
wavelength test packets (or slowly moving localized modes) toward the source.

\subsection{de~Broglie--Einstein ``harmony of phases''}

The classical de~Broglie picture emphasizes a relation between an internal ``clock'' (Compton-scale oscillation)
and an externally observed matter-wave phase.
In the present framework this relation is realized naturally: a localized soliton can support an internal carrier-scale
oscillation propagating at the microscopic wave speed $c$ in the effective substrate/branch metric, while the soliton
center-of-mass follows a timelike trajectory in the emergent kinematics.
The slow envelope phase then encodes the externally observed matter-wave behavior, while the internal carrier provides
the phase reference that enables Berry/Wilczek--Zee holonomy.

\subsection{Quantum phenomenology without fundamental probability}

By construction, the underlying dynamics is deterministic.
The probabilistic character of quantum measurements is attributed to the interaction of localized solitons with
wave packets that are constrained by Fourier localization and, in the lattice picture, by the finite Brillouin zone
(band-limited narrowband carriers cannot be localized arbitrarily sharply), together with threshold-like nonlinear exchange events in which small differences in phase and
amplitude can decide whether energy is transferred.
In this view, apparent point-like interactions reflect the structure of wave packets and nonlinear exchange,
not the existence of fundamental point particles.

\subsection{Limitations and open problems}

Several gaps remain before the theory fragment could be regarded as a closed replacement for established field theories:
\begin{itemize}
\item identifying constitutive laws and parameter regimes that simultaneously support (a) a stable ordered narrowband
sector and (b) robust localized solitons,
\item deriving the precise effective dynamics of curvature fluctuations of the Berry/Wilczek--Zee connection
\item quantifying the lattice-induced anisotropy/birefringence and establishing bounds (or suppression mechanisms)
that ensure effective Lorentz/isotropy constraints in accessible regimes,
\item formulating the dynamics and physical meaning of the inter-axis mixing factor $g_{\mathrm{mix}}$ (including when and how it
produces a protected near-degeneracy required for a non-Abelian Wilczek--Zee description),
(including whether a Maxwell/Yang--Mills-type term dominates in the relevant regimes),
\item clarifying the quantitative mapping between far-field holonomy labels and electromagnetic phenomenology,
\item understanding how multi-particle states and apparent entanglement-like correlations arise in a deterministic
wave medium.
\end{itemize}
The present goal is to place the theory on a consistent lattice-to-continuum footing so that these problems can be
attacked without ambiguity about the underlying substrate model.
